%==================================Problem Statement.tex=============================
\chapter{PROBLEM STATEMENT}
\label{ch:problem_statement}

This project addresses the critical need for affordable, deployable Security Information and Event Management (SIEM) capability in environments where dedicated SOC infrastructure is cost-prohibitive or operationally complex. The problem is framed through the following objectives and requirements.

\section{Problem Definition}

Small-to-medium organizations, educational institutions, and learning environments often lack the resources to implement enterprise-grade SIEM solutions. Commercial platforms such as IBM QRadar, Splunk, or Micro Focus ArcSight require significant licensing fees, specialized hardware, and dedicated personnel for deployment, operation, and maintenance. Open-source alternatives, while freely available, typically demand manual configuration, lack integrated correlation engines, and provide limited out-of-the-box visualization capabilities. Consequently, security teams in resource-constrained environments face fragmented visibility, manual correlation overhead, and delayed incident response times.

The SecuriSphere project was initiated to bridge this gap by providing a complete, lightweight SIEM platform deployable via Docker Compose on standard infrastructure. The platform implements a 3-layer pipeline model—collection, processing, and search—consistent with QRadar-style architectures while remaining accessible for educational, research, and small-scale operational use.

\section{Project Objectives}

This project aims to develop an intelligent, adaptable SIEM platform that utilizes rule-based detection and correlation engines to:

\begin{itemize}
    \item \textbf{Collect and Normalize Telemetry:} Implement a Linux-based agent for system log collection, metric extraction, and event transmission with HMAC-SHA256 integrity signing and offline SQLite buffering when network connectivity is unavailable.
    \item \textbf{Detect Security Threats:} Implement a rule-based detection engine with extensible checker registration supporting 14 built-in rule types, including failed login detection, brute force, privilege escalation, service integrity monitoring, disk/memory anomaly detection, file integrity surveillance, and network anomaly detection.
    \item \textbf{Correlate Related Events:} Implement a correlation engine with three matcher algorithms—sequence matching for ordered event patterns, co-occurrence matching for order-independent event sets, and cross-host matching for lateral movement detection—grouping related alerts into offenses with configurable clustering windows.
    \item \textbf{Reconstruct Attack Timelines:} Build an attack timeline engine that chains correlated events into chronological narratives, assigning confidence scores based on event type frequency, timeline compression, and privilege escalation indicators.
    \item \textbf{Provide UEBA Anomaly Detection:} Implement user and entity behavior analytics with rolling daily baselines, z-score thresholding, and severity classification, with configurable warm-up periods for baseline establishment.
    \item \textbf{Enable MITRE ATT\&CK Visualization:} Provide MITRE ATT\&CK tactic-coverage heatmaps and technique-level drill-down capabilities, linking detected events and alerts to their corresponding tactics and techniques for investigative efficiency.
    \item \textbf{Deliver a Complete SOC Dashboard:} Provide a web dashboard with real-time WebSocket updates, alert triage workflows, offense management, incident promotion, MITRE browsing, UEBA anomaly viewing, and investigation workspace integration.
    \item \textbf{Ensure Security and Reliability:} Implement JWT-based authentication with refresh token rotation, RBAC enforcement (admin/analyst/viewer roles), TOTP-based MFA, rate limiting, security headers, tamper-evident audit logging with SHA-256 hash chains, and session management with concurrent session limits.
\end{itemize}

\section{Requirements Analysis}

\subsection{Functional Requirements}

The platform shall provide the following functional capabilities:

\begin{longtable}{|l|l|}
\hline
\textbf{Requirement} & \textbf{Description} \\
\hline
RF-01 & Agent shall collect system auth logs, CPU/memory/disk metrics, and network flow data at configurable intervals (default: 10s for metrics, 30s for heartbeats) \\
\hline
RF-02 & Agent shall sign all event transmissions with HMAC-SHA256 using a per-agent API key and nonce validation \\
\hline
RF-03 & Agent shall buffer events locally in SQLite when network is down and auto-replay on reconnection \\
\hline
RF-04 & Backend shall validate HMAC signatures, nonce freshness (5 min), and timestamp skew (5 min) at the gatekeeper endpoint \\
\hline
RF-05 & Backend shall normalize event fields, map variant event types to canonical types, and store enriched events in PostgreSQL \\
\hline
RF-06 & Detection engine shall support extensible rule registration via an ABC-based checker registry pattern \\
\hline
RF-07 & Detection shall include 14 built-in rule types covering authentication, host health, system integrity, and network anomalies \\
\hline
RF-08 & Correlation engine shall evaluate sequence, co-occurrence, and cross-host patterns against recent event streams \\
\hline
RF-09 & Offense engine shall group related alerts into offenses using a configurable clustering window with sequential offense numbering \\
\hline
RF-10 & Attack timeline reconstruction shall chain correlated events into chronological narratives with MITRE technique assignment \\
\hline
RF-11 & UEBA shall compute rolling daily baselines for per-host and per-user metrics with z-score anomaly detection \\
\hline
RF-12 & MITRE ATT\&CK heatmap shall provide tactic-level coverage visualization with technique-level drill-down \\
\hline
RF-13 & Dashboard shall provide real-time WebSocket updates for alerts, offenses, host status, and live security feed \\
\hline
RF-14 & Authentication shall use JWT (HS256/RS256) in HttpOnly cookies with TOTP-based MFA support \\
\hline
RF-15 & RBAC shall enforce role-based permissions (admin, analyst, viewer) across all mutation endpoints \\
\hline
RF-16 & Audit log shall maintain tamper-evident SHA-256 hash chain integrity with advisory lock protection \\
\hline
RF-17 & Dashboard shall provide 29+ pages covering alerts, offenses, incidents, MITRE, UEBA, hosts, events, settings, and reports \\
\hline
\end{longtable}

\subsection{Non-Functional Requirements}

\begin{longtable}{|l|l|}
\hline
\textbf{Requirement} & \textbf{Description} \\
\hline
RNF-01 & System shall support deployment via Docker Compose on a single host with PostgreSQL and Redis \\
\hline
RNF-02 & Backend shall sustain sustained event ingestion of 100+ events/second per host \\
\hline
RNF-03 & WebSocket reconnection shall complete within 3 seconds of network disruption \\
\hline
RNF-04 & Alert rule configuration shall be manageable via API without restart \\
\hline
RNF-05 & UI shall support dark mode optimized for extended SOC analyst workflows \\
\hline
RNF-06 & Database schema migrations shall be managed via Alembic with versioned rollforward/rollback \\
\hline
RNF-07 & All code shall be covered by automated tests (backend: 74 test files, agent: 10 test files, frontend: 28 test files) \\
\hline
RNF-08 & Platform shall operate on a single Linux host with $\ge$2 CPU cores, 4GB RAM, and Docker \\
\hline
\end{longtable}

\section{Feasibility Study}

The feasibility of the SecuriSphere platform is supported by the following evidence:

\begin{itemize}
    \item The FastAPI framework provides high-performance async I/O for concurrent WebSocket and HTTP connections, validated by 74 backend test files covering detection, correlation, offense engine, UEBA, MITRE mapping, and API endpoints.
    \item The Next.js 14 React framework with TypeScript enables type-safe component development and SSR, with 28 frontend test files (20 unit tests + 8 E2E tests) ensuring UI reliability across dashboard pages, MITRE browsing, UEBA visualization, and offense management.
    \item The Python agent with SQLite offline buffer and HMAC signing has been tested across 10 test cases, confirming integrity and connectivity resilience including offline buffer overflow, heartbeat staleness, and HMAC verification failure scenarios.
    \item Docker Compose deployment configuration provides infrastructure-agnostic deployment for educational and small-scale operational environments.
    \item The rule registry pattern enables extensible detector addition without modifying core engine code, demonstrated by 14 registered rule types across authentication, host health, system integrity, and network anomaly categories.
    \item The platform supports sustained event ingestion of 100+ events/second per host with horizontal scalability via Redis queuing and PostgreSQL connection pooling.
\end{itemize}

The primary technical risk lies in the rule-based detection engine's limitation to known patterns; however, the modular checker registry design mitigates this by enabling future ML/model-based checker integration. The UEBA baseline computation requires a minimum of 3 days of aggregated stats before anomalies can fire, which is addressed via a configurable warm-up period and daily aggregation cron job.

\section{Aggregated Test Suite Summary}

\begin{longtable}{|l|l|l|}
\hline
\text{Component} & \text{Test Files} & \text{Test Functions (approx)} \\
\hline
Backend & 74 & ~500+ test functions across async unit, integration, and property-based suites \\
Agent & 10 & ~67 test functions covering collector, buffer, integrity, config, and sender modules \\
Frontend & 28 & ~20 unit test files + 8 E2E test files covering dashboard pages, MITRE browsing, UEBA, and offense workflows \\
\hline
\end{longtable}
\end{parameter>